\documentclass[10pt,a4paper]{article}
\usepackage[T1]{fontenc}
\usepackage{lmodern,amsmath,amssymb,amsthm,mathtools,microtype,geometry,enumitem,hyperref,booktabs}
\geometry{left=19mm,right=19mm,top=17mm,bottom=17mm}
\hypersetup{hidelinks,pdftitle={Probability Metrics and Couplings - Part 3 - Revised Study Notes}}
\setlength{\parindent}{0pt}\setlength{\parskip}{3.3pt}
\setlist[itemize]{topsep=4pt,itemsep=3pt}
\newtheorem{proposition}{Proposition}
\newtheorem{lemma}{Lemma}
\newcommand{\R}{\mathbb R}\newcommand{\E}{\mathbb E}\newcommand{\N}{\mathbb N}
\newcommand{\TV}{d_{\mathrm{TV}}}
\begin{document}
\begin{center}
{\Large\bfseries PART 3: Probability Metrics and Couplings}\par
{\normalsize Mathematical reading notes --- revised edition, 2026}
\end{center}
\vspace{8pt}
\textbf{Topics.} Metrics on spaces of probability laws; total variation distance (\S9.1); coupling and its relation to total variation (\S9.2); Wasserstein distance (\S12.4); quantile coupling on the real line; and an example comparing different optimal couplings.

\section*{1. Metric spaces and probability distributions}
Informally, a \emph{space} consists of a set together with a specified mathematical structure. A \emph{metric} on a set $A$ is a map $d:A\times A\to[0,\infty)$ satisfying for all $x,y,z\in A$
\begin{align*}
&d(x,y)=0\ \Longleftrightarrow\ x=y \quad\text{(identity of indiscernibles)},\\
&d(x,y)=d(y,x) \quad\text{(symmetry)},\\
&d(x,z)\le d(x,y)+d(y,z)\quad\text{(triangle inequality)}.
\end{align*}
Probability measures can themselves be regarded as points in a metric space, allowing us to compare whole distributions rather than individual sample outcomes.

\section*{2. Total variation distance (\S9.1)}
For probability measures $\mu$ and $\nu$ on the same measurable space $(S,\mathcal S)$, define
\[
\TV(\mu,\nu)=\sup_{A\in\mathcal S}|\mu(A)-\nu(A)|.
\]
Thus total variation records the largest difference in the probabilities assigned to the same event. The formula is a metric on the space of probability measures: symmetry and identity follow directly, and the triangle inequality follows from
\[
|\mu(A)-\lambda(A)|\le|\mu(A)-\nu(A)|+|\nu(A)-\lambda(A)|
\]
for every measurable $A$, followed by taking the supremum.

For example, if $x,y\in\R$, then $\TV(\delta_x,\delta_y)=0$ when $x=y$ and $1$ otherwise. In the latter case, the singleton event $\{x\}$ distinguishes the two point masses.

\subsection*{Two normalized densities: common and excess mass}
Suppose $f_1,f_2\ge0$ are densities with respect to one reference measure $\lambda$ and
\[
\int_S f_1\,d\lambda=\int_S f_2\,d\lambda=1.
\]
Set $h=f_1-f_2$, and write $h_+=\max\{h,0\}$, $h_-=\max\{-h,0\}$. As $\int h\,d\lambda=0$,
\begin{equation}\label{eq:mass}
\int_S h_+\,d\lambda=\int_S h_-\,d\lambda
=\frac12\int_S|h|\,d\lambda.
\end{equation}
The common mass is
\[
\int_S(f_1\wedge f_2)\,d\lambda
=1-\int_S(f_1-f_2)_+\,d\lambda
=1-\frac12\int_S|f_1-f_2|\,d\lambda,
\]
where $f_1\wedge f_2=\min\{f_1,f_2\}$. These identities formalize the intuition that the probability masses shared by the two laws cancel, while half their $L^1$ difference counts the excess mass on either side.

\begin{proposition}[Total variation of densities]
If $\mu$ and $\nu$ have densities $f_1$ and $f_2$ with respect to a common measure $\lambda$, then
\[
\TV(\mu,\nu)=\frac12\int_S|f_1-f_2|\,d\lambda.
\]
\end{proposition}
\begin{proof}
Let $A=\{x:f_1(x)>f_2(x)\}$. Then
\[
\mu(A)-\nu(A)=\int_A(f_1-f_2)\,d\lambda
=\int_S(f_1-f_2)_+\,d\lambda
=\tfrac12\int_S|f_1-f_2|\,d\lambda.
\]
Thus the supremum defining $\TV$ is at least the right side. Conversely, for every measurable $B$,
\[
\mu(B)-\nu(B)=\int_B(f_1-f_2)\,d\lambda
\le\int_S(f_1-f_2)_+\,d\lambda,
\]
and exchanging $f_1$ and $f_2$ gives the corresponding bound for $\nu(B)-\mu(B)$. Taking absolute values and then the supremum establishes equality.
\end{proof}
When $S$ is finite or countable and $p_1,p_2$ are probability mass functions, the formula becomes
\[
\TV(\mu,\nu)=\frac12\sum_{x\in S}|p_1(x)-p_2(x)|.
\]
The summand compares the \emph{two distributions at the same point $x$}, not values at two arbitrary points.

\subsection*{Examples}
For Bernoulli distributions,
\[
\TV\bigl(\mathrm{Ber}(p),\mathrm{Ber}(q)\bigr)
=\frac12\bigl(|(1-p)-(1-q)|+|p-q|\bigr)=|p-q|.
\]
Consider also distributions over three income groups (low, medium, high):
\[
\mu=(0.2,0.6,0.2),\qquad\nu=(0.5,0.2,0.3).
\]
Their total variation distance is
\[
\TV(\mu,\nu)=\tfrac12(|0.2-0.5|+|0.6-0.2|+|0.2-0.3|)=0.4.
\]
Here $0.4$ is the amount of mass that must be reassigned in an optimal discrete matching, consistent with the \emph{common-mass} identity above.


\subsection*{Detailed derivation of the common-mass identity}
For the normalized densities $f_1,f_2$, let $A=\{x:f_1(x)>f_2(x)\}$. Then $(f_1-f_2)_+=(f_1-f_2)\mathbf1_A$ and $(f_2-f_1)_+=(f_2-f_1)\mathbf1_{A^c}$. Since the two densities integrate to one,
\begin{align*}
0&=\int_S(f_1-f_2)\,d\lambda\\
 &=\int_A(f_1-f_2)\,d\lambda
   -\int_{A^c}(f_2-f_1)\,d\lambda.
\end{align*}
The two positive quantities on the last line must coincide. Their sum is $\int|f_1-f_2|\,d\lambda$, so each equals one half of that sum. Moreover,
\[
\int_S\min(f_1,f_2)\,d\lambda
=\int_A f_2\,d\lambda+\int_{A^c}f_1\,d\lambda
=1-\int_A(f_1-f_2)\,d\lambda.
\]
These are exactly the identities from the preceding lemma. This proof explains the ``common mass versus excess mass'' interpretation used later to construct couplings.

\section*{3. Coupling two distributions (\S9.2)}
A \emph{coupling} of probability measures $\mu_1$ and $\mu_2$ is a joint law $\pi$ on $S_1\times S_2$ whose marginals are $\mu_1$ and $\mu_2$. Equivalently, it is a pair of random variables $(X_1,X_2)$ defined on the same probability space with the prescribed marginal laws. The choice of coupling fixes their dependence; the independent coupling is just one possibility.

For example, given distributions of people's heights and weights, different couplings describe different possible assignments of heights to weights while leaving each marginal distribution unchanged. A coupling is therefore more information than knowing the two marginals separately.

\begin{proposition}[Coupling inequality for total variation]
For any coupling $(X_1,X_2)$ of probability laws on the same space for which equality of outcomes is measurable,
\[
\TV(\mu_1,\mu_2)\le P(X_1\ne X_2).
\]
On standard Borel spaces, a \emph{maximal coupling} attains equality.
\end{proposition}
\begin{proof}
For each measurable $A\subseteq S$,
\begin{align*}
|\mu_1(A)-\mu_2(A)|
&=\left|\E\bigl[\mathbf1_A(X_1)-\mathbf1_A(X_2)\bigr]\right|\\
&\le\E\bigl|\mathbf1_A(X_1)-\mathbf1_A(X_2)\bigr|\\
&\le P(X_1\ne X_2).
\end{align*}
The final inequality uses the fact that the two indicator functions coincide whenever $X_1=X_2$. Take the supremum over all $A$. The existence of a maximal coupling is a separate standard result, not proved here.
\end{proof}


\subsection*{How a maximal coupling uses common mass}
When two probability laws admit densities $f_1,f_2$ with respect to a common measure, let $c(x)=\min\{f_1(x),f_2(x)\}$ and $\alpha=\int c\,d\lambda$. We have $1-\alpha=\TV(\mu_1,\mu_2)$. If $0<\alpha<1$, the following explicit construction couples the laws:
\begin{enumerate}[leftmargin=*,label=(\roman*)]
\item With probability $\alpha$, sample $Z$ from density $c/\alpha$ and set $X_1=X_2=Z$.
\item With probability $1-\alpha$, sample $X_1$ and $X_2$ from the residual densities $(f_1-c)/(1-\alpha)$ and $(f_2-c)/(1-\alpha)$, respectively.
\end{enumerate}
The resulting marginals are $f_1$ and $f_2$. Since the residual densities have disjoint supports up to null sets, their independent draws disagree almost surely in the second branch. Thus
\[
P(X_1\ne X_2)=1-\alpha=\TV(\mu_1,\mu_2).
\]
The limiting cases $\alpha=0$ and $\alpha=1$ follow directly by using only the relevant branch. This argument explains the existence of a maximal coupling for the density case; the general standard-Borel statement can be established using analogous decomposition ideas.

\subsection*{Bernoulli example: independent versus maximal coupling}
Let $X_1\sim\mathrm{Ber}(p)$ and $X_2\sim\mathrm{Ber}(q)$ with $0\le p<q\le1$. If $X_1,X_2$ are independent,
\begin{align*}
P(X_1\ne X_2)
&=P(X_1=1,X_2=0)+P(X_1=0,X_2=1)\\
&=p(1-q)+(1-p)q=p+q-2pq.
\end{align*}
For $p=0.5$, $q=0.6$, this probability is $0.5$, giving a valid but loose upper bound on the actual total variation distance $0.1$.

To construct an optimal coupling, take independent variables $\theta_1\sim\mathrm{Ber}(r_1)$ and $\theta_2\sim\mathrm{Ber}(r_2)$, and set
\[
X_1=\theta_1,\qquad X_2=\theta_1\vee\theta_2.
\]
The marginal conditions require
\[
r_1=p,\qquad q=r_1+(1-r_1)r_2,
\qquad r_2=\frac{q-p}{1-p}
\]
(when $p<1$, which holds here). Then
\[
P(X_1\ne X_2)=P(\theta_1=0,\theta_2=1)
=(1-p)r_2=q-p=\TV\bigl(\mathrm{Ber}(p),\mathrm{Ber}(q)\bigr).
\]
This achieves the lowest possible disagreement probability. For the numerical example it equals $0.1$.


\subsection*{Joint probabilities for the Bernoulli coupling}
The independent coupling of two Bernoulli laws has the joint probability table
\[
\begin{array}{c|cc}
 & X_2=0 & X_2=1 \\\hline
X_1=0 & (1-p)(1-q) & (1-p)q\\
X_1=1 & p(1-q) & pq
\end{array}
\]
so the probability of a mismatch is the sum of the two off-diagonal entries. For $p=0.5$ and $q=0.6$ this is $0.2+0.3=0.5$, while the total variation distance is $|p-q|=0.1$.

In the maximal coupling with $p<q$, the joint table is instead
\[
\begin{array}{c|cc}
 & X_2=0 & X_2=1 \\\hline
X_1=0 & 1-q & q-p\\
X_1=1 & 0 & p
\end{array}
\]
which has the same marginals and mismatch probability $q-p$. The improvement comes from replacing unnecessary mismatches by shared outcomes. This concrete example is a useful reminder that knowing the marginals alone never specifies their dependence.

\section*{4. Wasserstein distance (\S12.4)}
For $p\ge1$, let $\mathcal P_p(\R)$ denote probability measures with finite $p$th absolute moments. The $p$-Wasserstein distance is
\[
W_p(\mu,\nu)=\left(\inf_{(X,Y)\in\Gamma(\mu,\nu)}
\E|X-Y|^p\right)^{1/p},\qquad\mu,\nu\in\mathcal P_p(\R).
\]
The infimum runs over all couplings of the two distributions. In this framework total variation also has a coupling representation, but uses the cost $\mathbf1_{\{X\ne Y\}}$, whereas Wasserstein distance uses $|X-Y|^p$. The cost for $W_2$ grows quadratically with displacement; the analogy with physical transport energy is a useful intuition, not an identity with mechanical energy $\tfrac12 mv^2$.

A standard result, used without proof here, is that $W_p$ is a metric on $\mathcal P_p(\R)$. Thus $W_p(\mu,\nu)=0$ exactly when the laws coincide. The $p$th root ensures the distance scales linearly under changes of spatial units.

\subsection*{Quantile functions and optimal coupling}
The quantile (generalized inverse CDF) is
\[
Q_\mu(u)=\inf\{x\in\R:F_\mu(x)\ge u\},\qquad0<u<1.
\]
For $U\sim\operatorname{Unif}(0,1)$ and $p\ge1$,
\begin{equation}\label{eq:quantile}
W_p(\mu,\nu)
=\left(\E\left|Q_\mu(U)-Q_\nu(U)\right|^p\right)^{1/p}
=\left(\int_0^1|Q_\mu(u)-Q_\nu(u)|^p\,du\right)^{1/p}.
\end{equation}
This is the one-dimensional \emph{quantile coupling}: equal quantile levels are paired. The absolute value is essential, in particular when $p$ is not an even integer. Optimality is a standard result from one-dimensional optimal transport, taken as given in these reading notes.

\subsection*{A three-city example: different metrics, different couplings}
Place cities $A,B,C$ at $-x,0,x$, where $x>0$. Consider the supply and demand laws
\[
\mu=\tfrac12\delta_{-x}+\tfrac12\delta_0,
\qquad\nu=\tfrac12\delta_0+\tfrac12\delta_x.
\]
Two possible couplings are:
\begin{center}
\begin{tabular}{@{}lll@{}}\toprule
Coupling & Pairs, each with mass $1/2$ & Interpretation\\ \midrule
$\gamma_0$ & $(-x,x)$ and $(0,0)$ & Maximize exact matches\\
$\gamma_1$ & $(-x,0)$ and $(0,x)$ & Pair equal quantile levels\\ \bottomrule
\end{tabular}
\end{center}
The mismatching probabilities and quadratic transport costs are
\begin{align*}
P_{\gamma_0}(X\ne Y)&=\tfrac12,&
\bigl(\E_{\gamma_0}|X-Y|^2\bigr)^{1/2}&=\sqrt2\,x,\\
P_{\gamma_1}(X\ne Y)&=1,&
\bigl(\E_{\gamma_1}|X-Y|^2\bigr)^{1/2}&=x.
\end{align*}
Because $\TV(\mu,\nu)=1/2$, the first coupling is maximal for total variation. By the quantile formula \eqref{eq:quantile}, the second coupling is optimal for $W_2$, giving $W_2(\mu,\nu)=x$. Thus a coupling optimal for one probability metric need not be optimal for another.


\subsection*{Transport plans for the three-city illustration}
The two couplings of $\mu=\frac12\delta_{-x}+\frac12\delta_0$ and $\nu=\frac12\delta_0+\frac12\delta_x$ can also be expressed as probability measures on pairs:
\begin{align*}
\gamma_0&=\tfrac12\delta_{(-x,x)}+\tfrac12\delta_{(0,0)},\\
\gamma_1&=\tfrac12\delta_{(-x,0)}+\tfrac12\delta_{(0,x)}.
\end{align*}
The first plan moves half the mass by distance $2x$ and leaves half unchanged. Its quadratic cost is therefore
\[
\E_{\gamma_0}|X-Y|^2
=\frac12(2x)^2+\frac12(0)^2=2x^2.
\]
The second plan moves each half-mass portion a distance $x$; its quadratic cost is
\[
\E_{\gamma_1}|X-Y|^2
=\frac12x^2+\frac12x^2=x^2.
\]
After taking square roots, these costs become $\sqrt2\,x$ and $x$, respectively, just as in the table above. By contrast, the first plan is better for the \emph{zero-one mismatch cost} because it preserves one half of the shared mass exactly. Thus the ``best'' way to couple a pair of laws depends on which notion of distance we intend to measure.

\subsection*{Why the quantile construction is well suited to one dimension}
The quantile construction pairs $Q_\mu(u)$ and $Q_\nu(u)$ at the same rank $u$. It is monotone: greater values of $u$ are paired with greater values in both distributions. For convex transport costs $|x-y|^p$, $p\ge1$, such monotone matching minimizes the expected cost on the real line. For a simple discrete example, if the first law puts equal mass at $a_1\le a_2$ and the second at $b_1\le b_2$, the noncrossing pairing $(a_1,b_1),(a_2,b_2)$ is no more expensive than the crossing pairing $(a_1,b_2),(a_2,b_1)$. The general quantile identity follows from this monotone-rearrangement principle.

\section*{Reading note}
The next part uses $W_2$ and independent optimal couplings to study the normalized sums in Tanaka's central limit theorem.
\end{document}
